\makeatletter
\def\input@path{{./}{../}{../../}{../../../}{../../../../}{preamble/}{../preamble/}{../../preamble/}{../../../preamble/}}
\makeatother

\input{preamble_exams.tex}
\usepackage{needspace}
\usepackage{titling}
\setlength{\droptitle}{-2em}
\pretitle{\par\begingroup\centering\LARGE}
\posttitle{\par\endgroup}
\preauthor{}\postauthor{}
\predate{}\postdate{}
\maketitlehookd{\vspace{-1em}}
\linespread{1.12}
\setlength{\parskip}{5pt}
\newgeometry{margin=0.65in,top=0.55in,bottom=0.65in}

\title{Grade 11 Term 2 Home Practice -- C10\\\large Investigation of merchant fee revenue\\[0.25em]
	\normalsize Global Economics -- 2026--2027}
\author{}
\date{}

\begin{document}
	\maketitle

	\setcounter{tocdepth}{1}
	\setcounter{secnumdepth}{2}
	\phantomsection
	\label{toc}
	\tableofcontents

	\section*{Evaluated criteria}
	\begin{enumerate}[label=\textbf{C\arabic*},start=10]
		\item Integrates research design, data analysis, and confidence-interval inference to investigate an economic or financial question, justifying the selected procedure and communicating estimates, precision, assumptions, limitations, and evidence-based conclusions.
	\end{enumerate}

	\textbf{How to practise.} Try each task before opening its blue solution link. The arrows return to the contents, problem list or current problem.

	\section{Worked Solution Examples}
	\label{sec:worked-solution-examples}

	\SectionProblemsTableOfContents{
		\SectionProblemLink{prob:c10-1}{Design the merchant revenue study (C10)} \\
		\SectionProblemLink{prob:c10-2}{Analyse the sampled merchant revenue (C10)} \\
		\SectionProblemLink{prob:c10-3}{Report the finding to the bank (C10)}
	}

	% ----------------------------------------------------------------
	\subsection{Design the merchant revenue study (C10)}
	\label{prob:c10-1}

	A bank wants to know whether the mean weekly net fee revenue per enrolled small-business merchant under a payments package exceeds a budgeting benchmark of 12.90 thousand USD. It has a complete roster of 500 merchants enrolled throughout a specified reference week. The controller can extract the week's net fees from the ledger, after rebates and reversals, for selected merchants. The bank wants a 95\% interval with a margin of error at most 0.50 thousand USD. A pilot suggests that the population SD may be near 0.90 thousand USD.

	\QuestionsTasks[sec:worked-solution-examples]{
		{10}/{sol:c10-1-a}/
		{Define the parameter and observational unit. Propose a way to select records and state one check on how the data are measured.},
		{10}/{sol:c10-1-b}/
		{Using $z^*=1.96$ and the pilot SD as a planning input, find an initial sample-size estimate. Explain how the analysis should verify the precision target after data collection.}
	}

	\clearpage

	\Needspace{5\baselineskip}
	\phantomsection
	\label{sol:c10-1-a}
	\subsubsection*{C10 -- Specify the question, population and design}

	Ask whether the mean net fee revenue per merchant in the reference week exceeds 12.90 thousand USD. The observational unit is one enrolled merchant and $\mu$ is the mean of the 500 merchant totals for that week. Number the roster and randomly select distinct IDs without replacement. Extract one weekly total per ID using the same rule for fees, rebates and reversals. Include eligible zero-fee merchants and verify unusual totals against the ledger. A fixed week gives a common measurement window.

	\secInicio[prob:c10-1][sec:worked-solution-examples]

	\Needspace{5\baselineskip}
	\phantomsection
	\label{sol:c10-1-b}
	\subsubsection*{C10 -- Plan the sample and explain the precision target}

	For a first approximation, require $1.96(0.90)/\sqrt n\le 0.50$, or $n\ge (1.96\cdot0.90/0.50)^2=12.446784$. This gives 13 as the minimum for the planning calculation. It is $13/500=2.6\%$ of the roster. The pilot SD is approximate, and a small-sample $t$ critical value will exceed 1.96. The team should budget for extra records and compare the realized $t^*s/\sqrt n$ with 0.50 before judging whether the final interval meets the precision target.

	\secInicio[prob:c10-1][sec:worked-solution-examples]

	\clearpage

	% ----------------------------------------------------------------
	\subsection{Analyse the sampled merchant revenue (C10)}
	\label{prob:c10-2}

	The bank randomly selects 16 distinct merchants from the 500-member roster before inspecting any of their revenues. All 16 records are available. The ledger extract gives each sampled merchant's total net fee revenue for the reference week, in thousand USD.
	\[
	12,\ 13,\ 14,\ 13,\ 12,\ 14,\ 13,\ 13,\ 14,\ 12,\ 13,\ 14,\ 13,\ 12,\ 14,\ 13
	\]
	For a two-sided 95\% interval, the Student's $t$ critical value at 15 degrees of freedom is 2.13145.

	\QuestionsTasks[sec:worked-solution-examples]{
		{10}/{sol:c10-2-a}/
		{Calculate $\bar x$ and $s$, describe the sample shape, check the design and choose a confidence-interval procedure. Use the Term 1 sample variance convention with divisor $n-1$.},
		{10}/{sol:c10-2-b}/
		{Construct the 95\% interval, report the achieved margin of error and compare it with the bank's target.}
	}

	\Needspace{5\baselineskip}
	\phantomsection
	\label{sol:c10-2-a}
	\subsubsection*{C10 -- Summarise the data and justify the interval}

	The counts at 12, 13 and 14 are 4, 7 and 5. The sample is roughly symmetric with no isolated extreme values. This supports the shape assumption, although 16 observations give limited information about the population's shape. The random sample is $16/500=3.2\%$ of the roster, supporting approximate independence.
	\[
	\sum x_i=209,\qquad \sum x_i^2=2739,\qquad
	\bar x=\frac{209}{16}=13.0625,\qquad
	s=\sqrt{\frac{2739-209^2/16}{15}}\approx0.77190.
	\]
	Both $\bar x$ and $s$ are in thousand USD. The response is quantitative and the population SD is unknown, so use a one-sample $t$ interval for $\mu$, with $df=15$.

	\secInicio[prob:c10-2][sec:worked-solution-examples]

	\Needspace{5\baselineskip}
	\phantomsection
	\label{sol:c10-2-b}
	\subsubsection*{C10 -- Estimate the population mean and check precision}

	\[
	\widehat{SE}=\frac{s}{\sqrt{16}}\approx0.19298,\qquad
	E=2.13145\,\widehat{SE}\approx0.41132\text{ thousand USD}.
	\]
	Using unrounded values, the interval is $[12.65118,13.47382]$, or about $[12.65,13.47]$ thousand USD per merchant. The margin of 0.41132 thousand USD is below the target of 0.50 thousand USD. The interval procedure has approximately 95\% long-run coverage under the sampling and shape assumptions.

	\secInicio[prob:c10-2][sec:worked-solution-examples]

	\clearpage

	% ----------------------------------------------------------------
	\subsection{Report the finding to the bank (C10)}
	\label{prob:c10-3}

	The budgeting team needs an evidence brief about the 12.90 thousand USD benchmark. Use the study design, data and interval from the first two problems.

	\QuestionsTasks[sec:worked-solution-examples]{
		{10}/{sol:c10-3-a}/
		{Write a brief with the estimated mean, the interval and achieved precision, the conclusion about the benchmark, and one limitation on what the bank can infer.},
		{10}/{sol:c10-3-b}/
		{The bank wants a conclusion about revenue over a quarter. Define a suitable measurement and explain how you would account for repeated records from the same merchant.}
	}

	\Needspace{5\baselineskip}
	\phantomsection
	\label{sol:c10-3-a}
	\subsubsection*{C10 -- Communicate the result and its limits}

	\noindent
	\begin{minipage}[t]{0.44\textwidth}
	\vspace{0pt}
	\centering
	\begin{tikzpicture}
	\begin{axis}[width=\linewidth,height=3.8cm,axis lines=left,
	xmin=12.45,xmax=13.65,ymin=0,ymax=1.8,
	xlabel={Mean fee revenue (thousand USD)},
	xtick={12.5,12.9,13.5},ytick=\empty]
	\addplot[black,dashed] coordinates {(12.9,0) (12.9,1.13)};
	\addplot[blue!65!black,thick,mark=|] coordinates {(12.65118,0.7) (13.47382,0.7)};
	\addplot[blue!65!black,only marks,mark=*] coordinates {(13.0625,0.7)};
	\node[anchor=north,font=\small] at (axis cs:12.9,1.72) {Benchmark};
	\end{axis}
	\end{tikzpicture}

	The segment is the 95\% interval. The dot is the sample mean.
	\end{minipage}
	\hfill
	\begin{minipage}[t]{0.52\textwidth}
	\vspace{0pt}
	The 16 randomly selected merchants averaged about 13.06 thousand USD in net fees. We are 95\% confident that the mean for all 500 enrolled merchants in the reference week was about 12.65 to 13.47 thousand USD. The margin of about 0.41 thousand USD meets the 0.50 target.

	The interval contains values both below and above 12.90, so the study leaves open whether the population mean exceeded the benchmark. The result applies to the specified week and enrolled cohort. It gives no estimate of the revenue that would have occurred under a different package.
	\end{minipage}

	\secInicio[prob:c10-3][sec:worked-solution-examples]

	\Needspace{5\baselineskip}
	\phantomsection
	\label{sol:c10-3-b}
	\subsubsection*{C10 -- Propose a study for a broader decision}

	Define the new population as merchants eligible under a specified quarter-long rule. Take one total per randomly selected merchant for that quarter. If management wants an average weekly figure, divide each quarter total by the number of weeks in the common window before estimating the mean. Treating each merchant-week as independent would overlook correlation within merchants. Use a benchmark in units consistent with the new measurement.

	\secInicio[prob:c10-3][sec:worked-solution-examples]


\end{document}
